% problem_id: S014-lurie-ultracategories-proposition-1-4-9
% source_id: S014 (Jacob Lurie, Ultracategories)
% source_file: lurie-conceptual.pdf
% statement_locator: printed p. 20
% proof_locator: printed p. [20, 21]
% representation: PDF; transcribed from rendered page images (never from the text layer)
% status: complete (statement + author proof verbatim + reason + locators)
%
%% problem_id: L4-lurie-ultracategories-prop-1.4.9
%% source_id: lurie-ultracategories (Jacob Lurie, "Ultracategories")
%% source_file: lurie-conceptual.pdf
%% statement_locator: p. 20, Proposition 1.4.9
%% proof_locator: pp. 20--21 ("Proof of Proposition 1.4.9" begins on printed page 20 and continues on
%%  printed page 21)
%% representation: PDF; transcribed from rendered page images (pdftoppm -r 150 -png for whole pages, and
%%  -r 300 / -r 600 -png crops for every symbol-level reading), NOT from the text layer
%% status: TRANSCRIPTION ONLY
%%
%% Transcription notes (every symbol below was checked against magnified crops of the page images; nothing
%% was taken from the PDF text layer, whose glyphs are corrupted: it renders $\in$ as "2", $\alpha$ as
%% "↵", etc.):
%%  * Pages transcribed, in page order: 20, 21. Since the proof of Proposition 1.4.9 occupies the foot of
%%    printed page 20 and the head of printed page 21, the two sentences of page-20 context that frame the
%%    proposition (the last sentence of the paragraph preceding it) are reproduced as well; nothing else of
%%    page 20 is reproduced.
%%  * Boxed cross-reference numbers of the source (hyperref boxes: 1.3.7, 1.3.8, 1.4.1, 1.4.8, 1.4.9,
%%    8.1.4, 2.1.1, 2.1.2, 2.2.2, A.4.5) are transcribed as plain references.
%%  * Places where the printed text deviates from what the mathematics would suggest are kept exactly as
%%    printed and flagged by a %% [NOTE: ...] line; nothing has been silently corrected. Four such places
%%    were found and are flagged in situ: the undefined index letter "I" in condition (∗) and in the colimit
%%    diagram of (∗); "give by" for "given by"; the doubled "that that"; and the diagram named "σ" although
%%    the display immediately above it is named "τ".
%%  * No passage of the supplied pages was illegible, so there are no [ILLEGIBLE: ...] markers.
%%  * The three display diagrams are reproduced as arrays (as elsewhere in this corpus); each preserves every
%%    node, arrow and label of the source diagram, but the geometric layout is approximated. Where the source
%%    draws an arrow as a slanted line the array uses \nearrow / \searrow in the corresponding cell.

\documentclass[11pt]{article}
\usepackage{amsmath}
\usepackage{amssymb}
\usepackage{amsthm}
\usepackage{mathrsfs}
\usepackage[margin=1in]{geometry}

\theoremstyle{plain}
\newtheorem{proposition}{Proposition}

\begin{document}

%% --- page 20 ---

Proposition 1.4.8 is an immediate consequence of the following more general (and more precise) assertion:

\renewcommand{\theproposition}{1.4.9}
\begin{proposition}
Let $\mathcal{M}^+$ and $\mathcal{N}^+$ be categories and let $\mathcal{M} \subseteq \mathcal{M}^+$ and
$\mathcal{N} \subseteq \mathcal{N}^+$ be full subcategories which admit ultraproducts in $\mathcal{M}^+$ and
$\mathcal{N}^+$, respectively. Let $F^+ : \mathcal{M}^+ \to \mathcal{N}^+$ be a functor which carries
$\mathcal{M}$ into $\mathcal{N}$ and satisfies the following additional condition:

\smallskip
\noindent$(\ast)$ \emph{For every collection of objects $\{M_s\}_{s \in S}$ of $\mathcal{M}$ and every
ultrafilter $\mu$ on $S$, the maps}
\[
F^+(q_\mu^{S_0}) : F^+\Big(\prod_{s \in I} M_s\Big) \to F^+\Big(\int_S M_s d\mu\Big)
\]
\emph{exhibit $F^+(\int_S M_s d\mu)$ as a colimit of the diagram $\{F^+(\prod_{s \in I} M_s)\}_{\mu(I)=1}$
in the category $\mathcal{N}^+$.}
\smallskip

%% [NOTE: in condition (∗) as printed, the product is indexed by "I" (twice: in the displayed map and in the
%%  subscript "μ(I)=1" of the colimit diagram), whereas the collection is indexed by "S" and the letter "I"
%%  is never introduced. Reproduced exactly as printed; not corrected.]
%% [NOTE: "I" in the two occurrences above is a capital i, not the numeral 1 and not a lower-case l; this was
%%  checked on a 600 dpi crop.]

Let $F = F^+\vert_{\mathcal{M}}$, which we regard as a functor from $\mathcal{M}$ to $\mathcal{N}$, and
regard $\mathcal{M}$ and $\mathcal{N}$ as equipped with the ultrastructures \emph{give} by
Proposition 1.3.7. Then:

%% [NOTE: the source prints "the ultrastructures give by Proposition 1.3.7"; the participle "given" appears to
%%  be intended. Reproduced exactly as printed; not corrected.]

\begin{enumerate}
\item[(a)] For every collection of objects $\{M_s\}_{s \in S}$ of $\mathcal{M}$ and every ultrafilter $\mu$
on $S$, there is a unique map $\sigma_\mu : F(\int_S M_s d\mu) \to \int_S F(M_s)d\mu$ having the property
that, for each subset $S_0 \subseteq S$ with $\mu(S_0) = 1$, the diagram
\[
\begin{array}{ccc}
F^+\big(\prod_{s \in S_0} M_s\big) & \longrightarrow & \prod_{s \in S_0} F(M_s)\\[12pt]
\Big\downarrow{\ F^+(q_\mu^{S_0})\ } & & \Big\downarrow{\ q_\mu^{S_0}\ }\\[12pt]
F\big(\int_S M_s d\mu\big) & \xrightarrow{\ \ \sigma_\mu\ \ } & \int_S F(M_s)d\mu
\end{array}
\]
%% [NOTE: the top horizontal arrow of this diagram is unlabelled in the source; reproduced unlabelled.]
commutes.
\item[(b)] The morphisms $\{\sigma_\mu\}$ of (a) determine a left ultrastructure on the functor $F$.
\item[(c)] If $F^+$ preserves small products, then $\{\sigma_\mu\}$ is an ultrastructure on $F$ (that is, each
of the maps $\sigma_\mu : F(\int_S M_s d\mu) \to \int_S F(M_s)d\mu$ is an isomorphism).
\end{enumerate}
\end{proposition}

\begin{proof}[Proof of Proposition 1.4.9]
Assertion $(a)$ follows immediately from $(\ast)$, and $(c)$ is clear. To prove $(b)$, we show that
\emph{that} the maps $\{\sigma_\mu\}$ satisfy condition (2) of Definition 1.4.1 (conditions (0) and (1) are
immediate from the definitions). Fix a collection of objects $\{M_t\}_{t \in T}$ of $\mathcal{M}$ indexed by a
set $T$, a collection of ultrafilters $\{\nu_s\}_{s \in S}$ on $T$ indexed by a set $S$, and an ultrafilter
$\mu$ on the set $S$. Set $\lambda = \int_S \nu_s d\mu$. We wish to show that

%% [NOTE: the source prints "we show that that the maps"; the word "that" is doubled as printed. Reproduced
%%  exactly as printed; not corrected.]

%% --- page 21 ---

the diagram $\tau$ :
\[
\begin{array}{ccccc}
\int_T F(M_t)d\lambda & \xrightarrow{\ \ \Delta_{\mu,\nu_\bullet}\ \ } & & & \int_S\big(\int_T F(M_t)d\nu_s\big)d\mu\\[16pt]
\Big\uparrow{\ \sigma_\lambda\ } & & & & \Big\uparrow{\ \int_S \sigma_{\nu_s}d\mu\ }\\[16pt]
& & \int_S F\big(\int_T M_t d\nu_s\big)d\mu & & \\[16pt]
& & \Big\uparrow{\ \sigma_\mu\ } & & \\[16pt]
F\big(\int_T M_t d\lambda\big) & \xrightarrow{\ \ F(\Delta_{\mu,\nu_\bullet})\ \ } & & & F\big(\int_S(\int_T M_t d\nu_s)d\mu\big)
\end{array}
\]
%% [NOTE: the source draws the four long arrows of this diagram as straight lines connecting the corners; the
%%  array above places the two vertical labels $\sigma_\lambda$ (left, spanning the whole height) and
%%  $\sigma_\mu$ (right, between $F(\int_S(\int_T M_t d\nu_s)d\mu)$ and $\int_S F(\int_T M_t d\nu_s)d\mu$) at
%%  their printed heights.]
commutes (in the category $\mathcal{N}$). Let $u, v : F(\int_T M_t d\lambda) \rightrightarrows
\int_S(\int_T F(M_t)d\nu_s)d\mu$ be the maps given by clockwise and counterclockwise composition around the
diagram \emph{$\sigma$}. To show that $u = v$, it will suffice (by virtue of $(\ast)$) to show that $u$ and
$v$ agree after precomposition with $F(q_\lambda^{T_0}) : F(\prod_{t \in T_0} M_t) \to F(\int_T M_t d\lambda)$,
for every subset $T_0 \subseteq T$ with $\lambda(T_0) = 1$. Set $S_0 = \{s \in S : \nu_s(T_0) = 1\}$, so that
$\mu(S_0) = 1$. We are then reduced to verifying the commutativity of the diagram

%% [NOTE: the display immediately above this sentence is named "τ", but the sentence names the diagram "σ".
%%  Reproduced exactly as printed; not corrected.]

\[
\begin{array}{ccccc}
\int_T F(M_t) & \xrightarrow{\ \ \Delta_{\mu,\nu_\bullet}\ \ } & & & \int_S\big(\int_T F(M_t)d\nu_s\big)d\mu\\[6pt]
\Big\uparrow{\ \sigma_\lambda\ } & \nearrow{\ q_\lambda^{T_0}} & & \nearrow{\ q_\mu^{S_0}} & \Big\uparrow{\ \int_S \sigma_{\nu_s}d\mu\ }\\[6pt]
& & \prod_{s \in S_0}\big(\int_T F(M_t)d\nu_s\big) & & \\[6pt]
& \nearrow{\ \{q_{\nu_s}^{T_0}\}_{s \in S_0}} & \Big\uparrow{\ \prod_{s \in S_0} \sigma_{\nu_s}\ } & & \\[6pt]
\prod_{t \in T_0} F(M_t) & & \prod_{s \in S_0} F\big(\int_T M_t d\nu_s\big) & \xrightarrow{\ \ q_\mu^{S_0}\ } & \int_S F\big(\int_T M_t d\nu_s\big)d\mu\\[6pt]
\Big\uparrow & \nearrow{\ F(q_\lambda^{T_0})} & \Big\uparrow & & \Big\uparrow{\ \sigma_\mu\ }\\[6pt]
F^+\big(\prod_{t \in T_0} M_t\big) & \xrightarrow{\ \ F(\{q_{\nu_s}^{T_0}\}_{s \in S})\ } & F^+\Big(\prod_{s \in S_0}\big(\int_T M_t d\nu_s\big)\Big) & & \\[6pt]
& \nearrow{\ F(q_\lambda^{T_0})} & & \searrow{\ F(q_\mu^{S_0})} & \\[6pt]
F\big(\int_T M_t d\lambda\big) & \xrightarrow{\ \ F(\Delta_{\mu,\nu_\bullet})\ \ } & & & F\Big(\int_S\big(\int_T M_t d\nu_s\big)d\mu\Big)
\end{array}
\]
%% [NOTE: the source displays this as one large commutative diagram; the array above preserves every node,
%%  arrow and label of that diagram, but its geometric layout is approximated. The vertical arrows
%%  $F^+(\prod_{t \in T_0} M_t) \to \prod_{t \in T_0} F(M_t)$ (left) and
%%  $F^+(\prod_{s \in S_0}(\int_T M_t d\nu_s)) \to \prod_{s \in S_0} F(\int_T M_t d\nu_s)$ (middle) are
%%  unlabelled in the source; reproduced unlabelled. The two copies of the label $F(q_\lambda^{T_0})$
%%  (rows 6 and 8 of the array) are one and the same slanted arrow
%%  $F(\int_T M_t d\lambda) \to F^+(\prod_{t \in T_0} M_t)$ in the source, whose label is placed on it; the
%%  array splits label and arrow across two rows and therefore repeats it.]
%% [NOTE: in this diagram the label on the arrow into $F^+(\prod_{s \in S_0}(\int_T M_t d\nu_s))$ is printed
%%  as $F(\{q_{\nu_s}^{T_0}\}_{s \in S})$, with index set $S$, while the target product is indexed by
%%  $S_0 \subseteq S$. Reproduced exactly as printed; not corrected.]

in the category $\mathcal{N}^+$. Note that the inner region of the diagram commutes by the construction of the
maps $\sigma_{\nu_s}$, the upper region commutes by the construction of the Fubini transformation for the
ultrastructure on $\mathcal{N}$, the lower region commutes by the construction of the Fubini transformation
for the ultrastructure on $\mathcal{M}$, the region on the left commutes by the construction of
$\sigma_\lambda$, the region on the upper right commutes by functoriality, and the region on the lower right
commutes by the construction of $\sigma_\mu$. \hfill$\square$
\end{proof}

\end{document}
